\documentclass{article}
\usepackage[utf8]{inputenc}
\usepackage[english]{babel}
\usepackage{caption}
\usepackage{float}
\usepackage{subcaption}
\usepackage[table,xcdraw]{xcolor}

\usepackage[lmargin = {2.7cm},
rmargin = {2.7cm},
tmargin = {2.2cm},
bmargin = {2.2cm}
]{geometry}

\usepackage{comment}
\usepackage[hidelinks]{hyperref}
\usepackage[capitalise]{cleveref}

% TexStudio build config line: txs:///bibtex | txs:///biber | txs:///compile | txs:///view
\usepackage[
backend=biber,
style=numeric,
sorting=nty
]{biblatex}
\addbibresource{paper.bib}

\usepackage{graphicx}
\graphicspath{ {./} }

\title{Maschinenlesbare Ontologien und -Vokabulare in RDF aus dem Bereich der industriellen Fertigung}
\author{\textbf{Konrad Abicht} \\ k.abicht@gmail.com}
\date{TODO Datum}

\begin{document}

\maketitle

\paragraph{{\large Zusammenfassung}}

\paragraph{Einleitung:} Aktuell gibt es keine geprüfte Liste mit maschinenlesbaren, öffentlich-zugänglichen RDF-Ontologien und -Vokabularen mit Bezug zum Themengebiet der industriellen Fertigung.
Ähnliche Arbeiten, wie IndustryPortal, sind Community-getrieben und besitzen keine Prüfung der Metadaten, wodurch es zu Widersprüchen zwischen den Angaben auf dem Portal und in den RDF-Daten kommt.
Eine solche Liste soll interessierten Personen und Institutionen helfen, sich schneller als bisher einen Überblick über passende Ontologien zu verschaffen.

\paragraph{Methodik:} Es wurde im Zeitraum vom 01.11.2023 - 28.03.2024 eine Internet- und Literaturrecherche durchgeführt und eine CSV-Datei mit allen thematisch passenden bzw. referenzierten Ontologien und Vokabularen erstellt.
Zu jeder Ontologie wurden eine Reihe von Metadaten erfasst (z.B. Titel, Kurzbeschreibung, Projektseite, Version, letzte Änderung und Lizenz).
Es wurden alle Ontologien ignoriert, zu denen es zwar wissenschaftliche Publikationen gibt, jedoch keine öffentlich zugänglichen RDF-Daten. Weitere Auswertungen sind am Ende der Publikation zu finden.
Es wurden Skripte eingesetzt, um semi-automatisch referenzierte Ontologien zu ermitteln.

\paragraph{Resultate:} Als Ergebnis wurden \textbf{217} Ontologien bzw. Vokabulare für den Themenbereich industrielle Fertigung gefunden. Nur \textbf{XXXXXX} Ontologien besitzen vollständige Metadaten, bei dem Rest konnte mindestens ein Metadatum nicht ermittelt werden.  Die Rohdaten (CSV-Datei, Skripte etc.) stehen unter einer freien Lizenz (CC-BY 4.0) für alle Menschen zur freien Verfügung.

\paragraph{Diskussion:} Die Untersuchung wurde lediglich über einen Zeitraum von 5 Monaten durchgeführt.
Weiterhin wurde der thematische Fokus industrielle Fertigung im weitesten Sinne betrachtet, wodurch möglicherweise auch thematisch nicht vollständig passende Ontologien einbezogen wurden.

\paragraph{Fazit:} Der wissenschaftliche Beitrag dieser Arbeit ist eine händisch geprüfte Liste von Ontologien, die unter einer freien Lizenz zugänglich ist. Daneben gibt es eine Auswertung zu den Rechercheergebnissen, wodurch Rückschlüsse zum Zustand der betrachteten Ontologien gezogen werden können. Aufgrund der frei zugänglichen Forschungsdaten können nachfolgende Untersuchungen leichter durchgeführt werden.

\newpage

\section{Einleitung}

Formale Ontologien und Vokabulare sind heutzutage weit verstreut im Internet.
Aus diesem Grund ist eine umfassende Recherche nach Ontologien für ein konkretes Themengebiet oft müßsam und zeitraubend.
Man muss neben Spezifikationsseiten und wissenschaftlichen Publikationen auch verschiedene Ontologie-Portale konsultieren.
Der durchschnittliche Domain Experte besitzt häufig nicht das Hintergrundwissen in den Formalen Ontologien und Semantic Web Technologien \cite{sanfilippo2019formal}.

Es fehlt an Brücken zwischen den Formalen Ontologien, OWL-Ontologien und den verschiedenen wissenschaftlichen Disziplinen.
Diese Arbeit soll eine solche Brücke für die Bereiche OWL-Ontologien und industrielle Fertigung beisteuern.
Im Rahmen einer Internet- und Literaturrecherche wurden alle auffindbaren, öffentlich zugänglichen\footnote{RDF-Dokument ist über eine URL abrufbar.} und in RDF vorliegenden Ontologien (und Vokabulare) aus dem Themengebiet der industriellen Fertigung gesammelt.
Das Ziel war die Erstellung einer Liste, bestehend aus maschinenlesbaren Ontologien, die eine automatisierte Verarbeitung des modellierten Wissens erlauben (z.B. Training einer KI oder OWL Reasoning).
Interessenten können somit thematisch passende Ontologien aus der Liste auswählen, sich die RDF-Daten herunterladen und z.B. mit eigenen Axiomen ergänzen.
Bis zu diesem Zeitpunkt waren solche Testbeds/Experimente nur sehr begrenzt und mit hohem technischen Aufwand möglich.

Die Arbeit ist wie folgt strukturiert: [...] Kapitel 2 enthält ... in Kapitel 3 folgt ... usw.

\subsection{Öffentlich verfügbare Forschungsdaten}

Alle erstellten Forschungsdaten stehen in einem Github-Repository für die Öffentlichkeit zur Verfügung:

\textbf{TODO:} Link einfügen

Die Inhalte unterliegen dabei den Bedingungen der Creative Commons BY 4.0.
Etwaige Rechte der Ontologieautoren, die Metadaten betreffend (z.B. Zusammenfassungen, Namen), bleiben unberührt.

\subsubsection{Graph Visualisierung}

Skript: generate-yuml-diagram-config

- Vereinheitlichung der Prefixe bei den RDF-Daten (z.B. dc vs dcterms vs dcelements oder prov vs. prov\_namespace)

\section{Fachlicher Hintergrund}

In diesem Kapitel werden die notwendige Vorkenntnise für diese Arbeit kurz zusammengefasst.

\subsection{Konzepte und Technologien des Semantic Web}

Das Resource Description Framework (RDF)\footnote{https://www.w3.org/RDF/} ist ein Model für Datenaustausch im Internet.
RDF erweitert die Link-basierte Struktur des Internets.
In diesem Zusammenhang spielen URLs, URIs und IRIs eine wichtige Rolle.
URL steht für Uniform Resource Locator und gibt eine Adresse einer Resource im Internet an.
Jede URL kann als Uniform Resource Identifier (URI) angesehen werden.
Man nutzt URIs zur Identifikation von abstrakten oder physischen Resourcen.
URIs dürfen nur aus ASCII-Zeichen bestehen, weshalb man irgendwann Internationalized Resource Identifiers (IRI) eingeführt hat.
Der Begriff Dereferenzierbarkeit spielt im Rahmen dieser Arbeit eine wichtige Rolle.
Eine URL (URI, IRI) wird dereferenzierbar genannt, wenn man sie aufrufen kann und eine valide Antwort erhält (z.B. RDF-Daten für eine gegebene Ontologie IRI).
Die Web Ontology Language (OWL) baut auf RDF auf und ist eine Spezifikation des W3C, um Ontologien formal beschrieben und verteilen zu können.
Im Rahmen dieser Arbeit wird sich auf Ontologien im Dateiformat RDF/Turtle bzw. RDF/XML fokussiert.

\subsection{Ontologien und Kontrollierte Vokabulare}

Im Rahmen dieser Arbeit lag der Fokus auf maschinenlesbaren RDF Ontologien und Vokabularen.
Maschinenlesbarkeit ist gegeben, wenn die Ontologie in Form einer Textdatei (in einer RDF-Notation) vorliegt und die Datei über eine URL abgerufen werden kann.
Die Ontologie muss zudem ein Vokabular zur Beschreibung des Fachgebietes bereitstellen und eine logische Theorie (z.B. in Form von Axiomen, Regeln, Hierarchien) über das Fachgebiet zugrunde legen, die auf das Vokabular zurückgreift.
In der Praxis sind die Ontologien bei diesem Punkt manchmal lückenhaft, weshalb die Theorie nur in Ansätzen erkennbar sein muss.

Diese Kriterien basieren größtenteils auf den Ausführungen von Fabian Neuhaus aus \cite{neuhaus2018ontology}, mit ein paar Ergänzungen.
\textbf{Im Rahmen dieser Arbeit wird der Begriff Ontologie auch für Vokabulare genutzt, außer es wird explizit darauf hingewiesen.}
Der Grund für diese Festlegung sind die Rechercheergebnisse, die zeigen, dass manche Autoren ihre eigenen Arbeiten ihr Arbeiten sowohl als Ontologie als auch Vokabular bezeichneten.
Zum Beispiel bezeichnet Martin Hepp GoodRelations als standardisiertes Vokabular für Produkte, Preise und Unternehmensdaten, nutzt jedoch Ontologie als Quasi-Synonym, siehe das folgende Zitat: "\textit{GoodRelations is a standardized vocabulary (also known as "schema", "data dictionary", or "ontology") for product, price, store, and company data that can [...]}"\footnote{Quelle: \url{https://www.heppnetz.de/ontologies/goodrelations/v1.html}}.

\subsubsection{Ontologiearten}

\begin{figure}[H]
    \includegraphics[scale=0.4]{ontology-types}
    \centering
    \caption{Ontology Types}
\end{figure}

Die Rechercheergebnisse zeigten, dass die Ontologie-Autoren ihre Arbeiten überwiegend nach der etablierten Vierteilung einordnen.
Zuerst kommen die \textbf{Top-Level Ontologien} (Synonyme: Upper Ontology, Foundational Ontology). Sie beinhaltet fachgebiets-unabhängige Inhalte (Begriffe, Relationen, Axiome etc.) zur Beschreibung eines Ausschnitts der Realität.
Zu den bekanntesten gehören Suggested Upper Merged Ontology (SUMO), Descriptive Ontology for Linguistic and Cognitive Engineering (DOLCE) und Basic Formal Ontology (BFO).
Es folgt die Ebene der \textbf{Mid-Level Ontologie}, welche die Inhalte der Top-Level Ontologie ergänzt bzw. verändert. Sie besitzen einen geringen thematischen Bezug und dienen als inhaltliche Brücke zwischen einer Top-Level Ontologie und ihr direkt untergeordneten Kern- bzw. Task-Ontologien.
Danach folgt die Ebene der \textbf{Kern Ontologie} (Synonyme: Core Ontology, Domain Ontology).
Ihr Inhalt bezieht sich stärker auf ein Fachgebiet oder sie ergänzt die Inhalte anderer Ontologien.
Unter den Kern-Ontologien folgt die \textbf{Task-Ontologie} (Synonyme: Application Ontology).
Sie nutzt in der Regel die Inhalte von Ontologien höherer Ebenen und stellt eigene Inhalte für einen konkreten Anwendungsfall bereit.

Als Nutzer der hier vorgestellten Ontologien sollte man wissen, auf welcher Ebene sich eine Ontologie befindet.
Das impliziert die inhaltliche Ausprägung und etwaige Abhängigkeiten, die bei der späteren Verarbeitung eine Rolle spielen.
So übernimmt eine Kern Ontologie die Theorie einer Top-Level Ontologie (inkl. aller Implikationen), was bei Nichtbeachtung später zu Widersprüchen mit eigenen Axiomen führen kann.

\subsubsection{Ontology Design Pattern (ODP)}

Ontology Design Pattern (ODP) sind kleine unabhängige Ontologien, die einen sehr begrenzten Themenbereich modellieren und dabei einen hohen Grad an Wiederverwendbarkeit und Kombinierbarkeit mit anderen ODPs und Modellierungsansätzen besitzen \cite{gangemi2009ontology}.
Es besteht keine Beschränkung für deren Aufbau, Struktur und Ausrichtung.
Die ODPs lassen sich auch in die vorgestellte Vierteilung der Ontologien entsprechend einordnen.

\subsection{Lizenzinformationen}

Zu jeder Ontologie wurde die zugehörige Lizenz recherchiert.
Sie bestimmt die rechtlichen Rahmenbedingungen unter denen eine Ontologie benutzt werden darf.
Die verwendeten Lizenzen reichen von Public Domain bis hin zu restriktiven individuellen Lizenzregelungen.
Im Folgenden werden alle Lizenzen aufgelistet, die während der Recherche gefunden wurden:

\begin{enumerate}
    \item Apache License 2.0\footnote{\url{https://www.apache.org/licenses/LICENSE-2.0.html}}
    \item BSD License (2- und 3-Clause)\footnote{\url{https://opensource.org/license/bsd-2-clause}, \url{https://opensource.org/license/bsd-3-clause}}
    \item Creative Commons Licenses (in verschiedenen Arten und Versionen)\footnote{\url{https://creativecommons.org/licenses/?lang=en}}
    \item GNU Public License (verschiedene Versionen)\footnote{\url{https://www.gnu.org/licenses/licenses.en.html}}
    \item MIT License\footnote{\url{https://opensource.org/license/mit}}
    \item OGC Document License Agreement\footnote{\url{https://www.ogc.org/about-ogc/policies/document-license-agreement/}}
    \item Public Domain Dedication License (PDDL)\footnote{\url{https://opendatacommons.org/licenses/pddl/}}
    \item W3C Document License (2023 version)\footnote{\url{https://www.w3.org/copyright/document-license-2023/}}
\end{enumerate}


\section{Verwandte Arbeiten}

Es konnten keine Vorarbeiten gefunden werden, bei denen eine geprüfte Liste an Ontologien aus dem Themenbereich der industriellen Fertigung bzw. der Industrie im allgemeinen erstellt wurde.
Jedoch gibt es ein paar Onlinedienste und wissenschaftliche Publikationen, die ähnliche Anstrengungen betreffen.

\subsection{IndustryPortal}

IndustryPortal \cite{amdouni2023industryportal} (\url{https://industryportal.enit.fr/ontologies}) ist eine Open Source Plattform für die Entwicklung, Publikation und Pflege von Ontologien aus dem Themenbereich Industrie, welche im Rahmen des OntoCommons-Projektes\footnote{\url{https://ontocommons.eu/}} entwickelt wurde.
Auf der Plattform können registrierte Benutzer u.a. neue Ontologien eintragen und bestehende Metadaten verändern\footnote{\url{https://industryportal.github.io/}}.
Es wurden keine Angaben gefunden, die auf eine Überprüfung solcher Einreichungen hinweisen.

Die Hauptunterschiede zwischen IndustryPortal und dieser Arbeit liegen in der Methodik und dem Datenumfang.
Es konnte mehrfach nachgewiesen werden, dass Metadaten in IndustryPortal in manchen Fällen unvollständig sind bzw. im Widerspruch zu den RDF-Daten stehen.
Als Beispiel ist die Ontologie mit dem Titel "Industry 4.0 Knowledge Graph"\footnote{\url{https://industryportal.enit.fr/ontologies/I40KG}} zu nennen.
Dieser Titel weicht von dem Titel ab, der in den RDF-Daten verwendet wird: Industrial IoT Architecture Ontology (dcterms:title als Property).

Trotz dieser Beobachtungen ist IndustryPortal eine wertvolle Quelle, denn sie stellen kostenfrei eine Infrastruktur bereit, die eine langfristige Verbesserung der Ontologien und zugehörigen Metadaten erlaubt.
Es enthält zum aktuellen Zeitpunkt\footnote{28.03.2024} 109 Ontologien für den Themenbereich Industrie im weitesten Sinne.
Im Rahmen dieser Arbeit wurden jedoch über 217 Ontologie allein im erweiterten Themenbereich industrielle Fertigung ermittelt.

\subsection{OntoCommons Ontology Catalogue}

Der OntoCommons Ontology Catalogue\footnote{\url{https://data.ontocommons.linkeddata.es/index}} wurde im Rahmen des OntoCommons-Projektes erstellt und wird scheinbar weiterhin gepflegt.
Jeder kann über ein Online-Formular\footnote{\url{https://ontocommons.eu/node/146}} Ontologien für die Liste vorschlagen.
Aktuell\footnote{28.03.2024} enthält die Liste 37 Ontologien, die sich fokussieren auf die Themenbereiche industrielle Fertigung, Materialien und Bauwesen.
Im Rahmen dieser Arbeit wurden jedoch über 217 Ontologie allein im weiteren Themenbereich industrielle Fertigung ermittelt.
Der Umfang und die Qualität der Metadaten je Ontologie schwanken ebenfalls.
Bei manchen Ontologien ist kein Link auf die RDF-Daten angegeben, womit sich die Frage stellt, warum diese überhaupt erwähnt werden\footnote{Beispiel "MPFQ Ontology (Material-Process-Function-Quality)":

\url{https://data.ontocommons.linkeddata.es/vocabulary/MpfqOntology(material-process-function-quality)}}.

\subsection{Wissenschaftliche Publikationen}

Aufgrund der hohen thematischen Übereinstimmung wurde die Publikation "Ontologies for Industry 4.0" \cite{kumar2019ontologies} aufgenommen.
Die Autoren geben zuerst eine kurze Einführung in die Themengebiete Industrie 4.0, Factory 4.0 und Smart Manufacturing und zählen danach relevante Herausforderungen in diesem Zusammenhang auf, z.B. Mensch-Maschine-Kommunikation oder Datenanalyse.
Später im Papier werden vier Ontologien vorgestellt:
\begin{enumerate}
    \item \textbf{CORA:} Core Ontology for Robotics and Automation
    \item \textbf{ROA:} The Ontology for Autonomous Robotics\footnote{In der zugehörigen Publikation\cite{olszewska2017ontology} kürzen die Autoren die Ontologie jedoch mit ORA ab.}
    \item \textbf{ORArch}: Ontology for Robotic Architecture
    \item \textbf{O4I4}: Ontology for Industry 4.0
\end{enumerate}

Leider gaben die Autoren keine Links auf die zugehörigen RDF-Daten an, wodurch die Ontologien eigentlich nicht im Scope dieser Arbeit sind.
Jedoch konnte im Rahmen der Recherche ein Link auf die RDF-Daten der CORA Ontologie ermittelt werden.

Die Autoren des Papiers "Where to Publish and Find Ontologies? A Survey of Ontology Libraries" \cite{d2012publish} geben einen Überblick über Ontologie-Bibliotheken (inkludiert auch Ontology Directory, Ontology Repository, Ontology Archive).
Diese Publikation ist relevant, weil sie weiterhin nutzbare Ontologie-Bibliotheken auflistet.
Neben Metadaten wie Titel, Lizenz und die neuste Version, werden häufig auch die zugehörigen RDF Daten bereitgestellt.
Die folgenden Portale sind weiterhin nutzbar:

\begin{enumerate}
    \item BioPortal (\url{https://bioportal.bioontology.org/}, Themengebiete: Biomedizin)
    \item OBO Foundry (\url{https://obofoundry.org/}, Themengebiete: Biologie und Biomedizin)
    \item oeGOV (\url{http://www.oegov.us/}, Themengebiet: e-Government)
    \item Ontology Lookup Service (\url{https://www.ebi.ac.uk/ols4}, Themengebiete: Biomedizin)
    \item Ontology Design Patterns (\url{http://ontologydesignpatterns.org/wiki/Main\_Page}, viele Themengebiete, siehe auch \url{http://ontologydesignpatterns.org/wiki/Community:Domain})
    \item ONKI ontology server (\url{https://onki.fi/en/}, verschiedene Themengebiete)
\end{enumerate}

Für diese Arbeit sind Ontologie-Bibliotheken, wie hier aufgelistet, unerlässlich.
Sie decken zwar nur ein begrenztes Themengebiet ab, bieten in der Regel aber leichte Zugänglichkeit (z.B. mit Suchfunktion und Listen).

Meine Rechercheergebnisse bestätigen die Beobachtungen der Autoren, nämlich das es oft unvollständige Angaben zur Weiterverwendung und Lizenzierung von Ontologien gibt.

\section{Methodik}

Es wurde eine Literatur- und Internetrecherche durchgeführt und die passenden Ontologien in einer CSV-Datei gesammelt.
Es wurde das CSV-Format gewählt, weil CSV-Dateien einfach zu parsen sind und in jeder gängigen Programmier- und Skriptsprache unterstützt werden.
Interessenten können die CSV-Datei automatisiert, maschinell auswerten und die per URL referenzierten RDF-Daten der Ontologien direkt einlesen.
Das erlaubt durchgängige Maschinenlesbarkeit und einen hohen Grad an Automatisierung.
Bei der gesamten Recherche waren nur deutsch- und englischsprachige Inhalte von Interesse.

\subsection{Forschungsfragen}

Die folgenden Forschungsfragen wurden bei der Erstellung der Liste zugrundegelegt:

\begin{enumerate}
    \item Welche Ontologien für das Fachgebiet industrielle Fertigung (im weitesten Sinne) gibt es?
    \item Welche dieser Ontologien werden aktiv betreut bzw. wann war die letzte dokumentierte Aktivität im Projekt?
    \item Welche lizenzrechtlichen Regelungen muss man beim Einsatz einer Ontologie beachten?
\end{enumerate}

\subsection{Thematische Abgrenzung}

Die industrielle Fertigung ist thematisch sehr breit und viele Themenbereiche wie Fertigungsverfahren, Mess- und Prüftechnik und Automatisierung spielen dabei eine wichtige Rolle.
Zusätzlich gibt es eine Reihe von Trends, welche häufig interdisziplinär aufgestellt sind und vollkommen neue Thwmengebiete mitbringen (z.B. Big Data beim Smart Manufacturing).
Bei dieser Ausgangslage eine thematisch Abgrenzung für diese Untersuchung zu erstellen war nicht einfach.
Dadurch, dass Ontologien sich jedoch häufig in Onlinediesten befinden, die eine Schlüsselwort-basierte Suche besitzen, wurde eine Liste von Schlüsselwörtern erarbeitet: defect, digital twin, industry, machine, manufacturing, product, sensor, supply chain.
Sie repräsentieren zentrale Konzepte der industriellen Fertigung sowie angrenzenden Themengebieten / Trends.
Es wurde hier ein Mittelweg zwischen thematischer Genauigkeit und Handhabbarkeit angestrebt.
Die meisten Onlinedienste besaßen an sich schon wenige Ontologien, weshalb die Suche mit weiteren Schlüsselwörtern oft nicht nötig war, um den gesamten Bestand zu sichten.

\subsection{Recherchequellen}

Die Ontologieentwicklung hat starke wissenschaftliche Wurzeln, daher stand die Literaturrecherche am Anfang, um die Publikationen zu finden, welche eine Ontologie und ihre Inhalte näher vorstellen.
Für die Internetsuche wurde \textbf{Google Search} verwendet.
\textbf{Google Scholar}\footnote{\url{https://scholar.google.com/}} kam bei der Recherche nach wissenschaftlichen Publikationen zum Einsatz.

\subsubsection{Dezidierte Onlinedienste}

Daneben wurde in weiteren Onlinediensten gesucht, wobei nicht alle unmittelbar etwas mit Ontologien zu tun haben.
Die folgende Liste enthält alle Onlinedienste, die im Rahmen der Recherche genutzt wurden:

\begin{enumerate}
    \item \textbf{AURORAL Ontologien} (\url{https://auroral.iot.linkeddata.es/index.html}) - Liste von Ontologien, die für das AURORAL Projekt (\url{https://www.auroral.eu/#/}) entwickelt wurden. Es enthält u.a. Ontologien zum Thema Ladestationen und Energiekonsum.
    \item Projektseite der \textbf{Basic Formal Ontology} (kurz BFO, \url{https://basic-formal-ontology.org/users.html}) - Enthält eine Liste von Ontologie-Projekten, welche die Top-Level Ontologie BFO nutzen.
    \item \textbf{BioPortal} - Ontologie-Portal mit über 1094 Ontologien, vordergründig aus dem biomedizinischen Themenbereich\footnote{Geprüft am 26.03.2024}.
    \item \textbf{Basic Register of Thesauri, Ontologies \& Classifications} (kurz BARTOC, \url{https://bartoc.org}) - Eine Webseite mit angebundener Suchmaschine, die Knowledge Organization Systems, Vokabulare und Ontologien auflistet. Es wurden ausschließlich OWL Ontologien verwendet. Einträge, die nicht in Deutsch oder Englisch vorlagen wurden automatisch übersetzt und geprüft.
    \item \textbf{EU Vocabularies} (\url{https://op.europa.eu/en/web/eu-vocabularies/controlled-vocabularies}) - Eine Webseite u.a. mit Ontologien und Vokabularen, bereitgestellt von der Europäischen Union (bzw. einer ihr unterstellten Institution).
    \item \textbf{Github} (\url{https://www.github.com}) - Ein Onlinedienst zur Softwareentwicklung, der jedoch sehr häufig als Ort für Ontologie-Projekte genutzt wird.
    \item \textbf{IndustryPortal} (\url{https://industryportal.enit.fr/}) - Ontologie-Portal mit über 109 Ontologien\footnote{Geprüft am 26.03.2024.} aus der Industrie und angrenzenden Themenbereichen.
    \item \textbf{Linked Open Vocabularies} (\url{https://lov.linkeddata.es/dataset/lov/}) - Ein kuratierter Katalog von Vokabularen und Ontologien zur Beschreibung von Daten im Internet.
    \item \textbf{OntoCommons Ontology Catalogue} (\url{https://data.ontocommons.linkeddata.es/index}) - Manuell kuratierte Liste mit Ontologien aus den Themenbereichen Industrie, Produktion, Materialwissenschaften, Bauwirtschaften und weiteren.
    \item \textbf{ShowVoc} (\url{https://showvoc.op.europa.eu/}) - Ein Portal mit einer Liste von Ontologien, Vokabularen und anderen. Es wurden nur die Ontologien ausgewertet.
\end{enumerate}

Es wurden die Suchergebnisse für jeden dieser Suchbegriffe näher untersucht, insofern sie von einer Ontologie handelten.
Für den Fall, dass es neben Ontologien und Vokabularen noch andere Typen gab, so wurden diese Typen ignoriert.

\subsubsection{Auswertung von Namespaces und owl:import}

In Ontologien kommen häufig Namensräume und owl:import-Anweisungen zum Einsatz.
Namensräume dienen dazu häufig genutzte URLs abzukürzen, um damit Platz zu sparen und die Lesbarkeit zu erhöhen.
Bei owl:import-Anwensungen signalisiert man, dass die Ontologie auf Inhalte einer referenzierten Ontologie angewiesen ist bzw. diese die Ontologie erweitern.
In der Praxis gibt es verschiedene Auslegungen und Implementierungen\footnote{Gute Ausführungen zum Thema hier: \url{https://protegewiki.stanford.edu/wiki/How_Owl_Imports_Work}}.
Im Rahmen dieser Untersuchung wird von der Annahme ausgegangen, dass Ontologie-Autoren mit der Nutzung von Namespaces und owl:import signalisieren, dass die Inhalte der referenzierten Ontologie auf irgendeine Weise inhaltlich relevant sind.

Aus diesem Grund wurden beide ebenfalls ausgewertet.
Referenzierte Ontologien wurden nach den gleichen Kriterien geprüft und entsprechend eingeordnet.

\subsection{Auswahlkriterien für eine Ontologie}

Eine Ontologie musste für die Aufnahme die folgenden Mindestkriterien erfüllen:

\begin{enumerate}
    \item Die Ontologie liegt als Textdatei\footnote{Das schließt auch dynamisch generierte Textdateien mit ein.} in einer RDF-Notation (RDF/XML oder RDF/Turtle) vor und kann über eine URL heruntergeladen werden.
    \item Die Ontologie stellt ein Vokabular zur Beschreibung des Fachgebietes bereit und es liegt eine logische Theorie (z.B. in Form von Axiomen, Regeln, Hierarchien) über das Fachgebiet zugrunde, die auf das Vokabular zurückgreift.
    \item Die Inhalte der Ontologie haben einen direkten thematischen Bezug zu den aufgelisteten Schlüsselwörtern bzw. den damit zusammenhängenden Themen.
    \item In der Textdatei ist mindestens eine Instanz von owl:Ontology zu finden oder mindestens eine Klasse bzw. Property definiert.
\end{enumerate}

Es wurden alle Ontologien (nebst zugehöriger wissenschaftlicher Publikationen) ignoriert, bei denen keine RDF-Daten vorlagen.
Ohne die zugehörigen RDF-Daten kann nicht sichergestellt werden, dass eine Ontologie vollständig vorliegt.
Davon abgeleitete Arbeiten könnten später zu unerwarteten Fehlern und Widersprüchen führen.

\subsection{Erfasste Metadaten zu einer Ontologie}

Bei jeder Ontologie wurden die folgenden Metadaten erfasst:

\begin{itemize}
    \item \textbf{Name der Ontologie} - Der Name der Ontologie, der in den RDF-Daten bzw. der zugehörigen Dokumentation zu finden ist.
    \item \textbf{Manufacturing Industry related} - In der Spalte steht "yes", wenn die Ontologie einen direkten thematischen Bezug hat. Alle Ontologien ohne thematischen Bezug, die jedoch referenziert wurden, erhalten in der Spalte "no".
    \item \textbf{Abbreviation} - Die Abkürzung des Namens der Ontologie, sofern vorhanden.
    \item \textbf{Short description} - Eine kurze, prägnante Beschreibung über den Inhalt der Ontologie.
    \item \textbf{Project page or publication} - Eine URL auf die Projektseite, falls vorhanden. Alternativ eine URL zu einer Publikation über die Ontologie oder zu einem Ontologie-Portal mit weiteren Angaben.
    \item \textbf{Ontologie URI} - URI zur Ontologie. Sie ist in der Regel global eindeutig.
    \item \textbf{Latest version} - Falls vorhanden, eine Angabe zur neusten Version der RDF-Daten. Die neuste Version und das Datum der neusten, dokumentierten Änderung hängen nicht zusammen und können vollkommen verschiedene zeitliche Bezüge haben.
    \item \textbf{Latest activity found} - Eine Datumsangabe der neusten und dokumentierten Änderung an den RDF-Daten. Die Änderung kann, muss aber nicht zeitlich mit der neusten Version zusammenfallen.
    \item \textbf{RDF/XML file}: Eine URL zur RDF/XML-Datei der Ontologie, falls vorhanden.
    \item \textbf{RDF/Turtle file}: Eine URL zur RDF/Turtle-Datei der Ontologie, falls vorhanden.
    \item \textbf{Download location}: Manche RDF-Dateien besitzen keinen statischen Download-Link oder sind Teil eines ZIP Archivs. Diese findet man dann hier.
    \item \textbf{Autors/creators} - Eine Liste von Namen der Autoren oder beteiligten Gruppen/Unternehmen. Sollte es keine Angaben zu Autoren geben, werden hier alle Beitragende aufgelistet.
    \item \textbf{License} - Falls vorhanden, eine Angabe zu der verwendeten Lizenz.
\end{itemize}

Eine Bewertung der Inhalte (z.B. Qualität der RDF-Daten oder genutzte Lizenzen) fand nicht statt.
Aufgrund des Umfangs wurde die Liste der Metadaten auf die notwendigsten Angaben beschränkt.
Alle Angaben in der zugehörigen CSV-Datei wurden in Englisch erstellt, um eine größtmögliche Zugänglichkeit zu ermöglichen.

Bei der Auswertung der Metadaten wurden zuerst die RDF-Daten geprüft. Waren diese unvollständig, wurde auf die Projektseite zurückgegriffen, falls vorhanden.
Waren beide Quelle erfolglos, wurde versucht, die Metadaten über den Onlinedienst zu beziehen, der Daten über die Ontologie bereitstellte.
Lagen am Ende zu einem Metadatum immer noch keine Angaben vor, so wurde an der Stelle "Information not available" in der CSV-Datei hinterlegt.

\subsection{Inhaltliche Nacharbeiten}

Es wurden Einträge in der ontologies.csv inhaltlich angepasst, um ein Mindestmaß an Verständlichkeit zu gewährleisten:

\begin{enumerate}
    \item Unvollständige oder fehlende Angaben wurden durch Angaben von Ontologie-Portalen ergänzt, sofern inhaltlich passend.
    \item Manche Ontologien besaßen keine Angaben zu den Autoren, lediglich Beitragende. In diesem Fall wurden alle Beitragenden in das Autors/creators Feld übernommen.
    \item Konnten keine Autoren und Beitragenden ermittelt werden (z.B. VDI3682), so wurde versucht den Klarnamen der Autoren über Git-Commits zu ermitteln.
    \item Die Versionen und Datumsangaben (der letzten Aktivität) wurden vereinheitlicht, um eine Vergleichbarkeit herzustellen.
\end{enumerate}

\subsection{Eigene Hilfsskripte}

Zur Erleichterung der eigenen Arbeit wurden verschiedene PHP-Skripte geschrieben.
Sie erlaubten eine semi-automatische Einarbeitung der Ontologien und nachträgliche Validerung der Metadaten.
Alle Skripte sind im zugehörigen Github-Repository näher dokumentiert.

\section{Rechercheergebnisse}

Im Folgenden werden die wichtigsten Erkentnisse ausgewertet.
Die Datenbasis waren die \textbf{217} Ontologien, die einen thematischen Bezug zur industriellen Fertigung haben.

\begin{figure}[H]
    \begin{minipage}[c]{0.4\linewidth}
        % https://quickchart.io/chart-maker/edit/zm-678e36f9-512b-4266-9f58-53cad550212e
        \includegraphics[scale=0.14]{pie-complete-incomplete-metadata}
        \caption{}
    \end{minipage}
    \hfill
    \begin{minipage}[c]{0.4\linewidth}
        % https://quickchart.io/chart-maker/edit/zm-e38ea0a6-e0c3-4429-a876-c0331ab34670
        \centering
        \includegraphics[scale=0.14]{pie-files-download-location}
        \caption{}
    \end{minipage}%
\end{figure}

Nur \textbf{50} von 217 Ontologien besitzen vollständige Metadaten.
Bei dem Rest (N=167) gibt es mindestens ein Metadatum, bei dem der Wert fehlt und es daher "Information not available" beinhaltet.

\textbf{212} Ontologien besitzen entweder eine per URL abrufbare RDF/XML- oder RDF/Turtle-Datei.
Nur bei 5 Ontologien liegen die RDF-Daten zwar vor, der Zugriff benötigt jedoch Zusatzaufwand (z.B. das Entpacken eines ZIP-Archives).
Es ist nicht nachvollziehbar, warum die Ontologie-Autoren diese Hürden für die Konsumenten ihrer Ontologie eingebaut haben.

\begin{figure}[H]
    \begin{minipage}[c]{0.4\linewidth}
        % https://quickchart.io/chart-maker/edit/zm-95b07247-b436-4bc9-ab1b-049a812e8590
        \includegraphics[scale=0.14]{pie-activity-37}
        \centering
        \caption{}
    \end{minipage}
    \hfill
    \begin{minipage}[c]{0.4\linewidth}
        % https://quickchart.io/chart-maker/edit/zm-2bfdd2ee-2a98-4eac-b926-7097227c8c98
        \centering
        \includegraphics[scale=0.14]{pie-dereferenceable-ontology-uri}
        \caption{}
    \end{minipage}%
\end{figure}

Bei \textbf{89} Ontologien konnte eine Projektaktivität (z.B. Git-Commit) innerhalb der letzten 3 Jahre festgestellt werden.
Eine Projektaktivität vor über 3 Jahren, aber innerhalb der letzten 7 Jahre konnte bei \textbf{55} Ontologien festgestellt werden.
Und bei \textbf{26} Ontologien konnte keine oder nur eine Projektaktivität vor über 7 Jahren festgestellt werden.

Ein Großteil der Ontologie, \textbf{178 Stück}, besitzen eine dereferenzierbare Ontology URI.
Zur Überprüfung der Dereferenzierbarkeit wurde ein HTTP-Request abgesendet und der Response ausgewertet.
Als dereferenzierbar wurde eine Ontology URI angesehen, wenn der Response keinen Fehler zurücklieferte.

\begin{figure}[H]
    \includegraphics[scale=0.14]{bars-license-usage}
    \centering
    \caption{Meist genutzte Lizenzen (min. 10 Stück)}
\end{figure}

\textbf{180} Ontologien nutzten eine freie bzw. Open Source Lizenz und lediglich \textbf{5} Ontologien benutzten eine eigene Lizenzregelung, die sich keiner bekannten Open Source Lizenz direkt zuordnen lies.
Bei \textbf{32} Ontologien konnten keine Angaben zur Lizenz gefunden werden.
Diese Feststellung stimmt positiv, weil es damit jedem ermöglicht eigene Arbeiten auf den Ontologien aufzubauen, ohne Lizenzkosten oder rechtliche Auseinandersetzungen zu fürchten.
Es muss an dieser Stelle jedoch auch erwähnt werden, dass es zwei Ontologien gibt, die nur für nicht-kommerzielle Zwecke zur Verfügung stehen (EEPSA Ontology, Industry 4.0 Knowledge Graph).

Abschließend noch zwei Erwähnungen.
Es wurden Ontologien gefunden, die nicht mehr in ihrer Form existierten, weil sie z.B. in eine andere Ontologie übergangen sind. So GoodRelations ging 2012 in schema.org über. Ein anderes Beispiel ist die Ontologie RealEstateCore: sie gab ihre OWL-Ontologie in Version 4 auf und verlagerte sich dann auf SHACL und Digital Twin Definition Language Serialisierungen\footnote{Siehe auch \url{https://dev.realestatecore.io/docs/DTDL-or-SHACL} und \url{https://doc.realestatecore.io/3.3/full.html}}.

Und die Product Types Ontology\footnote{\url{http://www.productontology.org/}} stellt lediglich einen RDF-Daten-Dump mit nur 1000 der am häufigsten abgefragten Klassen als Ontologie bereit.

\section{Diskussion}

Die vorgestellte Ontologie-Liste hat keinen Anspruch auf Vollständigkeit.
Die Untersuchung wurde über einen begrenzten Zeitraum von 5 Monaten durchgeführt.
Weiterhin wurde der thematische Fokus industrielle Fertigung im weitesten Sinne betrachtet, wodurch auch thematisch nicht vollständig passende Ontologien einbezogen sein könnten.
An dieser Stelle wäre eine Nachuntersuchung mit Domainexperten sinnvoll.
In den Skripten wurden RDF-Bibliotheken genutzt, welche nicht alle RDF-Dateien lesen konnten\footnote{Fehler beim XML-Parsing}.
In diesen Fällen wurden die Dateien manuell ausgewertet.

\section{Fazit und Ausblick}

Die Rechercheergebnisse zeigen ein vielfältiges Bild.
Für das Themengebiet industrielle Fertigung und angrenzende Themengebiete konnten viele passende Ontologien gefunden werden.
Es scheint ein über die letzten 10 Jahre beständiges Interesse an Ontologien in diesen Bereichen zu geben, was die Zeitpunkte der ermittelten letzten Aktivitäten zeigen.
Ebenso ist positiv, dass einige Ontologien eine dezidierte Spezifikationsseite besitzen, auf der man neben den Metadaten, u.a. auch eine Liste der Klassen und Properties sowie andere Inhalte findet.

Jedoch war es bei dem Großteil der Ontologie aufwendig an die gewünschten Metadaten zu gelangen.
Es ist bemerkenswert, wie wenig Aufmerksamkeit die Ontologie-Autoren auf die Metadaten zu legen scheinen.
Sind es doch die Metadaten, die zuerst bei Suchanfragen in Ontologie-Portalen und Suchmaschinen herangezogen werden.
Die Suche nach den Metadaten wurde erschwert durch kaputte Web-Links (Stichwort: link rot), was mich dazu zwang manuell über verschiedene Onlinedienste an weitere Daten über eine Ontologie zu gelangen.

Diese Arbeit stellt ein stabiles Fundament für weitere Untersuchungen zu Ontologien in der industriellen Fertigung dar. Zudem erlaubt sie interessierten Personen und Unternehmen, die sich im RDF-Umfeld noch nicht so gut auskennen, einfacher an die Ontologie zu gelangen, ohne sich mit technischem Kleinklein beschäftigen zu müssen.

\section{Erwähnungen}

Ich danke Paul-Robert Kästner für die Sichtung und Zuarbeit der ersten Ontologie-Einträge.

Die Arbeit wurde im Rahmen KI-Werk Projektes durchgeführt, welches vom Bundesministerium für Bildung und Forschung (BMBF) gefördert wurde (https://www.cbasynergy.net/cba/ki-werk.html).

\newpage

\section{Anhang}

Ontologie-Liste zum Thema industrielle Fertigung (stark gekürzt, vollständige Darstellung in Rohdaten):
\begin{table}[H]
    \resizebox{\textwidth}{!}{%
        \begin{tabular}{|l|l|l|}
            \hline
            & Name & Project page, publication or RDF-file \\ \hline
            1 & 3D Modeling Ontology & http://bdi.si.ehu.es/bdi/ontologies/ExtruOnt/docs/ \\ \hline
            2 & Additive Manufacturing and Maintenance Operations Ontology & https://github.com/LA3D/ammo \\ \hline
            3 & Additive Manufacturing Ontology & https://www.nist.gov/programs-projects/systems-integration-additive-manufacturing \\ \hline
            4 & Additive Manufacturing and Maintenance Operations Ontology & https://github.com/LA3D/ammo \\ \hline
            5 & AURORAL Adapters Ontology & https://github.com/oeg-upm/auroral-adapters-ontology \\ \hline
            6 & AURORAL Cell-Tower Ontology & https://github.com/oeg-upm/auroral-cellTower-ontology \\ \hline
            7 & AURORAL Energy Ontology & https://github.com/oeg-upm/auroral-energy-ontology \\ \hline
            8 & AURORAL Logistic Ontology & https://github.com/oeg-upm/auroral-shipmentBiomass-ontology \\ \hline
            9 & AURORAL Vehicle Charger Ontology & https://github.com/oeg-upm/auroral-VehicleCharger-ontology \\ \hline
            10 & Automotive Industry Ontology & https://iurianu.rocks/developer/ontology-for-automotive-industry/ \\ \hline
            11 & Battery Interface Ontology & https://www.big-map.eu/dissemination/battinfo \\ \hline
            12 & Bicycle Vocabulary & http://www.ebusiness-unibw.org/ontologies/opdm/bicycle.html \\ \hline
            13 & Brick & https://brickschema.org/ \\ \hline
            14 & Building Automation and Control Network Ontology & https://bacowl.sourceforge.net/intro.html \\ \hline
            15 & Building Circularity Assessment Ontology & https://github.com/linmor-sys/BCAO \\ \hline
            16 & Building Element Ontology & https://pi.pauwel.be/voc/buildingelement/index-en.html \\ \hline
            17 & Building Ontology & https://bimerr.iot.linkeddata.es/def/building/ \\ \hline
            18 & Building Product Ontology & https://www.projekt-scope.de/ontologies/bpo/ \\ \hline
            19 & Building Topology Ontology & https://github.com/w3c-lbd-cg/bot \\ \hline
            20 & Capability and Skills Ontology based on Industry Standards & https://github.com/CaSkade-Automation/CaSk \\ \hline
            21 & Capability and Skills Ontology based on Manufacturing & https://github.com/CaSkade-Automation/CaSkMan \\ \hline
            22 & Car HiFi Vocabulary & http://www.ebusiness-unibw.org/ontologies/opdm/carhifi.html \\ \hline
            23 & Car Options Ontology & https://lov.linkeddata.es/dataset/lov/vocabs/coo/versions/2010-10-12.n3 \\ \hline
            24 & Chair Vocabulary & http://www.ebusiness-unibw.org/ontologies/opdm/chair.html \\ \hline
            25 & Classification of Business Functions & https://showvoc.op.europa.eu/\#/datasets/Classification\_of\_Business\_Functions\_\%28CBF\_1.0\%29/metadata \\ \hline
            26 & Clothing Product Information Ontology & http://www.ebusiness-unibw.org/ontologies/cpi/ns \\ \hline
            27 & Coffee Machine Vocabulary & http://www.ebusiness-unibw.org/ontologies/opdm/coffeemachine.html \\ \hline
            28 & Collaborative Manufacturing Service Ontology & https://zenodo.org/records/3374505 \\ \hline
            29 & Components for ExtruOnt & https://www.semantic-web-journal.net/system/files/swj2217.pdf \\ \hline
            30 & Computer Vocabulary & http://www.ebusiness-unibw.org/ontologies/opdm/computer.html \\ \hline
            31 & Context Aware System Observation Ontology & https://irstea.github.io/caso/OnToology/ontology/caso.owl/documentation/index-en.html \\ \hline
            32 & Cooker and Oven Vocabulary & http://www.ebusiness-unibw.org/ontologies/opdm/cookeroven.html \\ \hline
            33 & core & https://rds.posccaesar.org/ontology/plm/ont/core/ \\ \hline
            34 & Core Ontology for Robotics and Automation & https://github.com/srfiorini/IEEE1872-owl \\ \hline
            35 & Core Ontology for Robotics and Automation (Bare) & https://github.com/srfiorini/IEEE1872-owl \\ \hline
            36 & CORAX & https://github.com/srfiorini/IEEE1872-owl \\ \hline
            37 & Crystallography Domain Ontology & https://github.com/emmo-repo/domain-crystallography \\ \hline
            38 & CSS Ontology & https://github.com/CaSkade-Automation/CSS \\ \hline
            39 & DefectOnt & https://github.com/AndreaMazzullo/DefectOnt/ \\ \hline
            40 & Digital Buildings Ontology & https://github.com/google/digitalbuildings/ \\ \hline
            41 & Digital Camera Vocabulary & http://www.ebusiness-unibw.org/ontologies/opdm/digitalcamera.html \\ \hline
            42 & Digital Construction Materials & https://data.industryportal.enit.fr/ontologies/DCMATERIALS/submissions/1/download?apikey=019adb70-1d64-41b7-8f6e-8f7e5eb54942 \\ \hline
            43 & Digital Receiver Vocabulary & http://www.ebusiness-unibw.org/ontologies/opdm/digitalreceiver.html \\ \hline
            44 & DIN EN 61360 Ontology-Design-Pattern & https://github.com/hsu-aut/IndustrialStandard-ODP-DINEN61360 \\ \hline
            45 & DIN EN 62264-2 Ontology-Design-Pattern & https://github.com/hsu-aut/IndustrialStandard-ODP-DINEN62264-2 \\ \hline
            46 & DIN 8580 Ontology-Design-Pattern & https://github.com/hsu-aut/IndustrialStandard-ODP-DIN8580 \\ \hline
            47 & Dishwasher Vocabulary & http://www.ebusiness-unibw.org/ontologies/opdm/dishwasher.html \\ \hline
            48 & Distribution Element Ontology & https://pi.pauwel.be/voc/distributionelement/index-en.html \\ \hline
            49 & DogOnt: Ontology Modeling for Intelligent Domotic Environments & https://iot-ontologies.github.io/dogont/ \\ \hline
            50 & Domain Mechanical Testing & https://github.com/emmo-repo/domain-mechanical-testing \\ \hline
            51 & Domain Mechanical Testing Chemistry & https://github.com/emmo-repo/domain-mechanical-testing \\ \hline
            52 & Dryer Machine Vocabulary & http://www.ebusiness-unibw.org/ontologies/opdm/dryermachine.html \\ \hline
            53 & DVD Player and Blu-ray Player Vocabulary & http://www.ebusiness-unibw.org/ontologies/opdm/blurayplayer.html \\ \hline
            54 & eClass Products and Services Ontology & http://www.heppnetz.de/projects/eclassowl/ \\ \hline
            55 & EEPSA Ontology & https://iesnaola.github.io/eepsa/EEPSA/index-en.html \\ \hline
            56 & Elemental Multiperspective Material Ontology & https://emmc.eu/emmo \\ \hline
            57 & Elemental Multiperspective Material middle-level ontology & https://github.com/emmo-repo/EMMO \\ \hline
            58 & equipment & https://rds.posccaesar.org/ontology/plm/ont/equipment/ \\ \hline
            59 & ERA Vocabulary & https://showvoc.op.europa.eu/\#/datasets/ERA\_vocabulary/metadata \\ \hline
            60 & European Waste Classification for Statistics & https://showvoc.op.europa.eu/\#/datasets/ESTAT\_European\_Waste\_Classification\_for\_Statistics\_\%28EWC-Stat\_Rev.4\%29/metadata \\ \hline
            61 & Extruder Ontology & http://bdi.si.ehu.es/bdi/ontologies/ExtruOnt/docs/ \\ \hline
            62 & Extruder's sensors ontology & http://bdi.si.ehu.es/bdi/ontologies/ExtruOnt/docs/ \\ \hline
            63 & Facility Ontology & https://github.com/oeg-upm/cogito-facility-ontology \\ \hline
            64 & Fridge and Freezer Vocabulary & http://www.ebusiness-unibw.org/ontologies/opdm/refrigerator.html \\ \hline
            65 & Furniture Sector Ontology & https://industryportal.enit.fr/ontologies/FUNSTEP \\ \hline
            66 & Game Console Vocabulary & http://www.ebusiness-unibw.org/ontologies/opdm/gameconsole.html \\ \hline
            67 & Garment Vocabulary & http://www.ebusiness-unibw.org/ontologies/opdm/garment.html \\ \hline
            68 & General Process Ontology & https://gitlab.cc-asp.fraunhofer.de/ISC-Public/ISC-Digital/ontology/gpo \\ \hline
            69 & GRACE Ontology & https://industryportal.enit.fr/ontologies/GRACE \\ \hline
            70 & Grid2Onto & https://industryportal.enit.fr/ontologies/GRID2ONTO \\ \hline
            71 & Hair Dryer Vocabulary & http://www.ebusiness-unibw.org/ontologies/opdm/hairdryer.html \\ \hline
            72 & Home Hifi Vocabulary & http://www.ebusiness-unibw.org/ontologies/opdm/homehifi.html \\ \hline
            73 & IEEE Standard for Autonomous Robotics Ontology & https://github.com/hsu-aut/IndustrialStandard-ODP-IEEE1872-2 \\ \hline
            74 & IFC2X3 & https://github.com/buildingsmart-community/ifcOWL \\ \hline
            75 & IFC2X3\_TC1 & https://github.com/buildingsmart-community/ifcOWL \\ \hline
            76 & IFC4 & https://github.com/buildingsmart-community/ifcOWL \\ \hline
            77 & IFC4\_ADD1 & https://github.com/buildingsmart-community/ifcOWL \\ \hline
            78 & IFC4\_ADD2 & https://github.com/buildingsmart-community/ifcOWL \\ \hline
            79 & Industrial Maintenance Ontology & https://industryportal.enit.fr/ontologies/IMAMO \\ \hline
            80 & IOF Core Ontology & https://github.com/iofoundry/ontology/ \\ \hline
            81 & International System of Quantities & https://github.com/emmo-repo/EMMO \\ \hline
            82 & Industrial IoT Architecture Ontology & https://github.com/i40-Tools/StandardsOntology \\ \hline
            83 & Industrial Ontologies Foundry Supply Chain Reference Ontology & https://spec.industrialontologies.org/iof/ontology/supplychain/SupplyChainReferenceOntology/ \\ \hline
            84 & Industry 40 Knowledge Graph & https://industryportal.enit.fr/ontologies/I40KG \\ \hline
            85 & IoT Ontology & https://github.com/oeg-upm/cogito-iot-ontology \\ \hline
            86 & ISO 22400-2 Ontology-Design-Pattern & https://github.com/hsu-aut/IndustrialStandard-ODP-ISO22400-2 \\ \hline
            87 & Key Performance Indicator ontology & https://bimerr.iot.linkeddata.es/def/key-performance-indicator/ \\ \hline
            88 & Landline Phones Vocabulary & http://www.ebusiness-unibw.org/ontologies/opdm/landlinephone.html \\ \hline
            89 & M3-lite Taxonomy & https://github.com/fiesta-iot/ontology \\ \hline
            90 & M3 Ontology & http://sensormeasurement.appspot.com/m3\# \\ \hline
            91 & Maintenance Activity Ontology & https://github.com/uwasystemhealth/Paper\_Archive\_Maintenance\_Activity \\ \hline
            92 & Maintenance Activity Ontology & https://industryportal.enit.fr/ontologies/MNT-ACT \\ \hline
            93 & Maintenance Reference Ontology & https://spec.industrialontologies.org/iof/ontology/maintenance/Maintenance/ \\ \hline
            94 & Manufacturing & https://github.com/emmo-repo/EMMO \\ \hline
            95 & Manufacturing Semantics Ontology & https://www.academia.edu/download/30806306/Lemaignan2006.pdf \\ \hline
            96 & Manufacturing System Ontology & https://github.com/enegri/OFM \\ \hline
        \end{tabular}%
    }
\end{table}


% Please add the following required packages to your document preamble:
% \usepackage{graphicx}
\begin{table}[H]
    \resizebox{\textwidth}{!}{%
        \begin{tabular}{|l|l|l|}
            \hline
            & Name & Project page, publication or RDF-file \\ \hline
            97 & ManuService Ontology & https://industryportal.enit.fr/ontologies/MANUSERVICE \\ \hline
            98 & Materials & https://github.com/emmo-repo/EMMO \\ \hline
            99 & Materials Design Ontology & https://github.com/LiUSemWeb/Materials-Design-Ontology \\ \hline
            100 & Material properties ontology & https://bimerr.iot.linkeddata.es/def/material-properties/ \\ \hline
            101 & Material Science and Engineering Ontology & https://matportal.org/ontologies/MSEO \\ \hline
            102 & Mechanical Testing Ontology (MTO) & https://industryportal.enit.fr/ontologies/MTO \\ \hline
            103 & Microstructure domain ontology & https://github.com/emmo-repo/domain-microstructure \\ \hline
            104 & Microwave Vocabulary & https://ebusiness-unibw.org/ontologies/opdm/microwave.html \\ \hline
            105 & Mobile Phone Vocabulary & http://www.ebusiness-unibw.org/ontologies/opdm/mobilephone.html \\ \hline
            106 & Modem Vocabulary & http://www.ebusiness-unibw.org/ontologies/opdm/modem.html \\ \hline
            107 & Molecules And Materials Basic Ontology & https://github.com/daimoners/MAMBO \\ \hline
            108 & MPS500 & https://github.com/hsu-aut/MPS500-Capabilities \\ \hline
            109 & MPS500 AssembleCylinder & https://github.com/hsu-aut/MPS500-Capabilities \\ \hline
            110 & MPS500 AssembleThermometer & https://github.com/hsu-aut/MPS500-Capabilities \\ \hline
            111 & MPS500 CameraModule & https://github.com/hsu-aut/MPS500-Capabilities \\ \hline
            112 & MPS500 DrillingModule & https://github.com/hsu-aut/MPS500-Capabilities \\ \hline
            113 & MPS500 PropertyTypes & https://github.com/hsu-aut/MPS500-Capabilities \\ \hline
            114 & MPS500 RawCylinderSupplyModule & https://github.com/hsu-aut/MPS500-Capabilities \\ \hline
            115 & MPS500 RawThermometerSupplyModule & https://github.com/hsu-aut/MPS500-Capabilities \\ \hline
            116 & MPS500 ShippingModule & https://github.com/hsu-aut/MPS500-Capabilities \\ \hline
            117 & MPS500 StorageModule PutInStorage & https://github.com/hsu-aut/MPS500-Capabilities \\ \hline
            118 & MPS500 StorageModule RetrieveFromStorage & https://github.com/hsu-aut/MPS500-Capabilities \\ \hline
            119 & MPS500 Transport Module & https://github.com/hsu-aut/MPS500-Capabilities \\ \hline
            120 & MSDL (Manufacturing Service Description Language) & https://industryportal.enit.fr/ontologies/MSDL \\ \hline
            121 & Navigation Device Vocabulary & http://www.ebusiness-unibw.org/ontologies/opdm/navigation.html \\ \hline
            122 & Occupancy Profile ontology & https://github.com/oeg-upm/bimerr-occupant-behavior \\ \hline
            123 & oneM2M Base Ontology & https://git.onem2m.org/MAS/BaseOntology \\ \hline
            124 & Ontology for Maintenance Procedure Documentation (OMPD) Conditional Maintenance Task Ontology & https://industryportal.enit.fr/ontologies/OMPD-CMTO \\ \hline
            125 & Ontology for the Battery Value Chain & https://gitlab.cc-asp.fraunhofer.de/ISC-Public/ISC-Digital/ontology/bvco \\ \hline
            126 & Ontology model for Web of Things & http://iot.linkeddata.es/def/wot/index-en.html \\ \hline
            127 & Ontology of units of Measure & http://bdi.si.ehu.es/bdi/ontologies/ExtruOnt/docs/ \\ \hline
            128 & OntoSensor Device Ontology & https://mmisw.org/ont/univmemphis/sensor \\ \hline
            129 & OPC UA Core ontology & https://github.com/OntoUA/ua-nodeset-core-ont \\ \hline
            130 & OPC UA NodeSet ontology & https://github.com/OntoUA/ua-nodeset-core-ont \\ \hline
            131 & OPC UA Ontology-Design-Pattern & https://github.com/hsu-aut/IndustrialStandard-ODP-OPC-UA \\ \hline
            132 & OpenADR ontology & https://albaizq.github.io/OpenADRontology/OnToology/ontology/openADRontology.owl/documentation/index-en.html \\ \hline
            133 & OpenLink Product Features Ontology & http://www.openlinksw.com/ontology/features\# \\ \hline
            134 & Open Energy Ontology & https://openenergy-platform.org/ontology/ \\ \hline
            135 & PackML StateMachine Ontology-Design-Pattern & https://github.com/hsu-aut/IndustrialStandard-ODP-PackML \\ \hline
            136 & Paper Vocabulary & https://ebusiness-unibw.org/ontologies/opdm/paper.html \\ \hline
            137 & PCA 'Part 14' upper ontology & https://rds.posccaesar.org/ontology/lis14/ont/core/ \\ \hline
            138 & Platform Ontology & https://github.com/oeg-upm/cogito-platform-ontology \\ \hline
            139 & Portable Media Player Vocabulary & http://www.ebusiness-unibw.org/ontologies/opdm/portablemp.html \\ \hline
            140 & Position Ontology & https://github.com/srfiorini/IEEE1872-owl \\ \hline
            141 & Printer Vocabulary & https://ebusiness-unibw.org/ontologies/opdm/printer.html \\ \hline
            142 & process & https://rds.posccaesar.org/ontology/plm/ont/process/ \\ \hline
            143 & Process Ontology & https://github.com/oeg-upm/cogito-construction-process-ontology \\ \hline
            144 & Product Ontology & https://github.com/mvegetti/PRONTO/ \\ \hline
            145 & Product Types Ontology & http://www.productontology.org/ \\ \hline
            146 & Product Vocabulary & https://ns.inria.fr/provoc/v1/provoc\_v1.html \\ \hline
            147 & PSS Ontology & https://industryportal.enit.fr/ontologies/PSS \\ \hline
            148 & RealEstateCore Full & https://github.com/RealEstateCore/rec \\ \hline
            149 & Reference Generalized Ontological Model & https://github.com/MuhammadYahta/rgom \\ \hline
            150 & Reified Requirements Ontology & https://data.dnv.com/ontology/requirement-ontology/core/req-ont.html \\ \hline
            151 & Resistance Spot Welding Ontology & https://github.com/nsai-uio/RSWO \\ \hline
            152 & RESPOND Ontology & https://respond-project.github.io/RESPOND-Ontology/respond/index-en.html \\ \hline
            153 & RFID System Configuration Ontology & https://github.com/eleniTsalapati/ONTOLOGIES \\ \hline
            154 & RIVA InfoModel & https://github.com/hsu-aut/RIVA\_InfoModel \\ \hline
            155 & Resource, Material, Process, Function and Quality (rmpfq) ontology & https://github.com/zhengxiaochen/rmpfq\_ontology \\ \hline
            156 & ROMAIN: Reference Ontology for industrial Maintenance & https://industryportal.enit.fr/ontologies/ROMAIN \\ \hline
            157 & RPARTS & https://github.com/srfiorini/IEEE1872-owl \\ \hline
            158 & Safety Ontology & https://github.com/mahsa-teimourikia/Safety-Ontology \\ \hline
            159 & Safety Ontology & https://github.com/oeg-upm/cogito-safety-ontology \\ \hline
            160 & SAREF Ontology & https://saref.etsi.org/core/ \\ \hline
            161 & SAREF extension for the Automotive domain & https://saref.etsi.org/saref4auto/ \\ \hline
            162 & SAREF extension for building & https://saref.etsi.org/saref4bldg/ \\ \hline
            163 & SAREF extension for the Energy domain & https://saref.etsi.org/saref4ener/ \\ \hline
            164 & SAREF extension for the electric grid domain & https://saref.etsi.org/saref4grid/ \\ \hline
            165 & SAREF extension for the industry and manufacturing domain & https://saref.etsi.org/saref4inma/ \\ \hline
            166 & SAREF extension for the smart lifts domain & https://saref.etsi.org/saref4lift/ \\ \hline
            167 & SAREF4SYST: an extension of SAREF for typology of systems and their inter-connections & https://saref.etsi.org/saref4syst/ \\ \hline
            168 & SAREF4WATR is an extension of SAREF for Water & https://github.com/oeg-upm/S4WATR \\ \hline
            169 & SCOPRO (Supply Chain Process Ontology) & https://openreview.net/pdf?id=ricE3prhfH8 \\ \hline
            170 & SCORVoc & https://github.com/vocol/scor \\ \hline
            171 & SEAS Architecture ontology & https://ci.mines-stetienne.fr/seas/ArchitectureOntology \\ \hline
            172 & SEAS Building Ontology & https://ci.mines-stetienne.fr/seas/BuildingOntology \\ \hline
            173 & SEAS Device ontology & https://ci.mines-stetienne.fr/seas/DeviceOntology \\ \hline
            174 & SEAS Electric Light Source Ontology & https://ci.mines-stetienne.fr/seas/ElectricLightSourceOntology \\ \hline
            175 & SEAS Electric Power System Ontology & https://ci.mines-stetienne.fr/seas/ElectricPowerSystemOntology \\ \hline
            176 & SEAS Electric Street Light System Ontology & https://ci.mines-stetienne.fr/seas/StreetLightSystemOntology \\ \hline
            177 & SEAS Electric Vehicle ontology & https://ci.mines-stetienne.fr/seas/ElectricVehicleOntology \\ \hline
            178 & SEAS Photovoltaic ontology & https://ci.mines-stetienne.fr/seas/PhotovoltaicOntology \\ \hline
            179 & SEAS Smart Meter ontology & https://ci.mines-stetienne.fr/seas/SmartMeterOntology \\ \hline
            180 & SEAS Thermodynamic System ontology & https://ci.mines-stetienne.fr/seas/ThermodynamicSystemOntology \\ \hline
            181 & Semantic Sensor Network Ontology & https://www.w3.org/TR/vocab-ssn/ \\ \hline
            182 & Semantically Integrated Planning Model & https://industryportal.enit.fr/ontologies/SIMPM \\ \hline
            183 & Sensor Data ontology & https://bimerr.iot.linkeddata.es/def/sensor-data/ \\ \hline
            184 & Sensor, Observation, Sample and Actuator & https://www.w3.org/TR/vocab-ssn/ \\ \hline
            185 & Sharework Ontology for Human-Robot Collaboration & https://industryportal.enit.fr/ontologies/SOHO \\ \hline
            186 & Shaver Vocabulary & http://www.ebusiness-unibw.org/ontologies/opdm/shaver.html \\ \hline
            187 & Shoe Vocabulary & http://www.ebusiness-unibw.org/ontologies/opdm/shoe.html \\ \hline
            188 & Shredder Vocabulary & http://www.ebusiness-unibw.org/ontologies/opdm/shredder.html \\ \hline
            189 & SmartHomeWeather & https://paul.staroch.name/en/ \\ \hline
            190 & SmartProducts EADS & https://projects.kmi.open.ac.uk/smartproducts/ontology.html \\ \hline
            191 & SmartProducts Generic model & https://projects.kmi.open.ac.uk/smartproducts/ontology.html \\ \hline
            192 & SmartProducts Product model & https://projects.kmi.open.ac.uk/smartproducts/ontology.html \\ \hline
            193 & Statistical classification of products by activity & https://op.europa.eu/en/web/eu-vocabularies/dataset/-/resource?uri=http://publications.europa.eu/resource/dataset/cpa21 \\ \hline
            194 & Storage Media Vocabulary & http://www.ebusiness-unibw.org/ontologies/opdm/storagemedia.html \\ \hline
            195 & Tablet PC Vocabulary & http://www.ebusiness-unibw.org/ontologies/opdm/tablet.html \\ \hline
            196 & Television Vocabulary & https://ebusiness-unibw.org/ontologies/opdm/television.html \\ \hline
            197 & Top Level Ontology of Ontology-based-InformationFlow-Industry-40 & https://github.com/ko3n1g/Ontology-based-InformationFlow-Industry-4.0 \\ \hline
            198 & Units of measure (uom) & https://rds.posccaesar.org/ontology/plm/ont/uom/ \\ \hline
            199 & Utility vocabulary of OPDM Category Scheme based on the taxonomy of product types defined by Google & http://www.ebusiness-unibw.org/ontologies/opdm/category/google.owl \\ \hline
            200 & Vacuum Cleaner Vocabulary & http://www.ebusiness-unibw.org/ontologies/opdm/vacuum.html \\ \hline
            201 & VDI 2206 Ontology-Design-Pattern & https://github.com/hsu-aut/IndustrialStandard-ODP-VDI2206 \\ \hline
            202 & VDI 2860 Ontology-Design-Pattern & https://github.com/hsu-aut/IndustrialStandard-ODP-VDI2860 \\ \hline
            203 & VDI 3682 Ontology-Design-Pattern & https://github.com/hsu-aut/IndustrialStandard-ODP-VDI3682 \\ \hline
            204 & VDI 5100 Ontology-Design-Pattern & https://github.com/hsu-aut/IndustrialStandard-ODP-VDI5100 \\ \hline
            205 & VDI VDE NAMUR 2658 & https://github.com/hsu-aut/IndustrialStandard-ODP-VDIVDENAMUR2658 \\ \hline
            206 & Vehicle Sales Ontology for Semantic Web-based E-Commerce & https://www.heppnetz.de/ontologies/vso/ns \\ \hline
            207 & Versioning Ontology & https://data.ontocommons.linkeddata.es/vocabulary/VersioningOntology(veronto) \\ \hline
            208 & Vicinity core model & http://iot.linkeddata.es/def/core/index-en.html \\ \hline
            209 & Video Camera Vocabulary & http://www.ebusiness-unibw.org/ontologies/opdm/videocamera.html \\ \hline
            210 & Video Projectors Vocabulary & http://www.ebusiness-unibw.org/ontologies/opdm/videoprojector.html \\ \hline
            211 & Virtual Asset Representation Ontology & https://fundaciontekniker.github.io/a4blue-ontology/var/index-en.html \\ \hline
            212 & visualization4ExtruOnt & http://bdi.si.ehu.es/bdi/ontologies/ExtruOnt/docs/ \\ \hline
            213 & Volkswagen Vehicles Ontology & https://www.w3.org/2001/sw/sweo/public/UseCases/Volkswagen/Volkswagen.pdf \\ \hline
            214 & Washer Machine Vocabulary & http://www.ebusiness-unibw.org/ontologies/opdm/washingmachine.html \\ \hline
            215 & Waste Categories & https://showvoc.op.europa.eu/\#/datasets/ESTAT\_Waste\_categories/metadata \\ \hline
            216 & WoT Digital Twin Ontology & https://oeg-upm.github.io/WoT-DT-ontology/OnToology/dt-ontology.ttl/documentation/index-en.html \\ \hline
            217 & Z-BRE4K semantic model & https://www.z-bre4k.eu/ \\ \hline
        \end{tabular}%
    }
\end{table}






\medskip

\printbibliography

\end{document}

